% -----------------------------------------------------------------------
% Compiles as-is with pdflatex. No exotic packages, so it drops straight
% into Overleaf.
%
% Numbers marked \pending{} are the ones still to be filled from the runs
% listed at the end of each section. Everything else is measured.
% -----------------------------------------------------------------------
\documentclass[11pt,a4paper]{article}
\usepackage[margin=2.4cm]{geometry}
\usepackage{amsmath,amssymb}
\usepackage{graphicx}
\usepackage{booktabs}
\usepackage{xcolor}
\usepackage[colorlinks=true,linkcolor=blue,urlcolor=blue]{hyperref}
\usepackage{parskip}

\newcommand{\pending}[1]{\textcolor{red}{[#1]}}
\newcommand{\kv}{\ensuremath{\mathrm{s\,km^{-1}}}}
\newcommand{\taueff}{\ensuremath{\tau_{\rm eff}}}
\newcommand{\Pod}{\ensuremath{P_{\rm 1D}}}
\newcommand{\code}[1]{\texttt{\small #1}}

\title{Where the Lyman-$\alpha$ suppression in the FCT box comes from}
\author{P.~Colazo, N.~Padilla, F.~Stasyszyn}
\date{\today}

\begin{document}
\maketitle

\begin{abstract}
The 1D flux power spectrum of the Lyman-$\alpha$ forest is suppressed by
about $19\%$ on large scales in our FCT/PBH box relative to CDM. This note
is the record of everything we did to make that suppression go away. It
does not go away, but it is not what it first looked like. The suppression
is not in the matter field, where the two runs agree to $0.1\%$; it is in
the gas field, where they differ by $19\%$. The quick-Lyman-$\alpha$ star
formation scheme converts half the baryons in the FCT box against a seventh
in CDM, and the gas that survives traces the same matter with a lower bias.
The forest therefore does constrain this model, but through the gas bias
rather than through the primordial spectrum, and the amplitude of the
effect is set by a subgrid prescription without feedback. Every figure and
number here can be regenerated with the scripts at
\url{https://github.com/patricio-c/Lyman-Alpha-PBHs}.
\end{abstract}

\section{The question}

Both boxes are $40\,{\rm Mpc}/h$ on a side with $1024^3$ gas particles,
started from \emph{identical} initial conditions and evolved with SWIFT
using the quick-Lyman-$\alpha$ (QLA) prescription. The only difference is
the primordial power spectrum: CDM in one, the FCT broken spectrum
($n_b = 2$, $k_t = 10\,{\rm Mpc}^{-1}$) plus the Poisson isocurvature term
from PBH discreteness in the other. At $z=3$ we shoot 1536 sightlines of
2048 pixels each through both boxes, at the same transverse positions, and
compare the flux power spectra.

The ratio $\Pod^{\rm FCT}/\Pod^{\rm CDM}$ sits near $0.78$ at
$k = 0.003$~\kv{} and rises through unity at high $k$. The question this
note answers is whether that number is a property of the model or of how we
measured it.

We take the second possibility seriously enough to spend most of the note
on it, because the ratio is the whole result. If it survives, it is a
constraint on the model. If it does not, we have been reporting a bug.

\section{The two boxes}

\begin{table}[h]
\centering
\begin{tabular}{lrr}
\toprule
 & CDM & FCT \\
\midrule
box                       & \multicolumn{2}{c}{$40\,{\rm Mpc}/h = 58.7372\,{\rm Mpc}$} \\
gas particles in the ICs  & 134{,}217{,}728 & 134{,}217{,}728 \\
gas particles at $z=3$    & 115{,}318{,}346 & 67{,}341{,}600 \\
converted by QLA          & $14.1\%$ & $49.8\%$ \\
$m_{\rm gas}$             & \multicolumn{2}{c}{$9.4445\times10^{6}\,M_\odot$} \\
\bottomrule
\end{tabular}
\caption{The two runs. The initial conditions are the same file; the boxes
diverge only through the input power spectrum. The gas particle mass agrees
to $0.01\%$ with the \code{40-512-cc} run of Bolton et al.~(2017), which is
useful because their Figure~A4 gives us an external expectation for the
resolution trend.}
\end{table}

At $z=3$: $H(z) = 306.3\,{\rm km\,s^{-1}\,Mpc^{-1}}$, so the box spans
$4498\,{\rm km\,s^{-1}}$ and a pixel is $2.1963\,{\rm km\,s^{-1}}$. The
fundamental mode is $k = 0.0014$~\kv{} and the DESI DR1 window runs from
$0.001$ to $0.0318$~\kv{}. To convert, $k[{\rm Mpc}^{-1}] = 76.58\,k[\kv]$
and $k[h/{\rm Mpc}] = 112.45\,k[\kv]$. The FCT break at
$k_t = 10\,{\rm Mpc}^{-1}$ sits at $0.131$~\kv{}, well above the DESI
window, which is worth keeping in mind: the suppression we are chasing is
nowhere near the scale where the model was designed to act.

We analyse only below half the Nyquist frequency. Above that the estimator
aliases and the ratio between two runs develops structure that is pure
grid.

\section{The measurement, with and without normalisation}
\label{sec:t0}

The first objection to deal with, and the one raised after the talk, is the
normalisation. The two runs do not produce the same effective optical
depth. We rescale both to the observed value,
$\taueff = 0.3719$ at $z=3$ (Turner et al.~2024), through a multiplicative
factor $A$ on $\tau$:

\begin{equation}
F = e^{-A\tau}, \qquad
-\ln\langle F\rangle = \taueff^{\rm obs} .
\end{equation}

This is a nonlinear operation on the flux, so in principle it can move
power between scales, and one could reasonably suspect it of manufacturing
the ratio. We find $A_{\rm CDM} = 0.78961$ and $A_{\rm FCT} = 0.87198$, so
the two runs are being pushed by different amounts and the concern is not
idle.

Figure~\ref{fig:t0} settles it. The flux PDF and the power spectrum are
shown in both states: raw, with $A=1$ in both runs, and rescaled to a
common \taueff. The suppression is present in the raw fields, before
anything is normalised. Rescaling changes the ratio at $k=0.003$~\kv{} by
\pending{run \code{tests/t0\_rescaling.py}} and in the direction that makes
it slightly deeper, not shallower.

\begin{figure}[h]
\centering
\includegraphics[width=\textwidth]{../figures/t0_rescaling.png}
\caption{The rescaling is not the cause. Top: flux PDF with no rescaling
and after forcing both runs to the observed \taueff. Middle: the power
spectrum in the same two states. Bottom left: the ratio computed both ways
on the same axes; bottom right: the difference between them.}
\label{fig:t0}
\end{figure}

One number from the same test deserves its own line, because the first
version of the talk got it backwards. After rescaling, the FCT box has
about $21$ times \emph{more} fully transparent pixels ($F > 0.99$) than
CDM, not fewer. That is consistent with everything below, and it is the
first hint of where this is going.

\section{Five things that do not explain it}

Each of these was run to completion. None of them is worth a figure, so
they get a paragraph each.

\textbf{Sampling variance.} A paired bootstrap over sightlines and a
split-half test give a suppression significant at $20\sigma$. With 1536
matched sightlines through identical initial conditions, the pairing
removes most of the cosmic variance, and what is left is far too small.

\textbf{The deposition scheme.} We recomputed without the SPH kernel
normalisation. The ratio survives unchanged. This also disposes of an
earlier version of the analysis in which the deposition lacked an impact
parameter and overestimated $\tau$ by a factor $5.6$; that bug is fixed and
the ratio is not sensitive to the fix.

\textbf{Dense gas removed by QLA.} Cutting all gas above $\Delta = 100$ and
above $\Delta = 30$ from both runs changes the ratio by less than $1\%$.
Whatever carries the effect is not the densest surviving gas.

\textbf{Temperature.} Imposing the same $T$--$\rho$ relation on both runs
($\gamma = 1.40$) accounts for about $2\%$. The thermal histories do differ
--- they must, since the collapse histories differ --- but not enough.

\textbf{Saturation.} Masking saturated pixels at increasing aggressiveness
does not move the ratio monotonically; it wanders between $0.81$ and $0.85$
with no trend. Saturated systems in the sense of Bolton et al.~(2017) are
not the mechanism.

We also considered shot noise from the smaller number of FCT gas particles
and dropped it without running anything: the stochastic term falls by a
factor $\sim 6000$ between $k=0.002$ and $k=0.1$~\kv{}, while the effect we
see is flat over that range. A white-noise contaminant cannot produce a
flat ratio.

\section{The deficit is in the gas, not in the matter}
\label{sec:3d}

At this point we stopped testing the pipeline and measured the 3D power
spectra directly, which is what we should have done first.

\begin{table}[h]
\centering
\begin{tabular}{lrrr}
\toprule
$k\,[{\rm Mpc}^{-1}]$ & matter & baryons & gas only \\
\midrule
0.107 & 1.000 & 1.000 & 0.808 \\
0.249 & 1.000 & 1.000 & 0.807 \\
0.608 & 1.003 & 1.003 & 0.824 \\
1.020 & 1.011 & 1.011 & 0.866 \\
2.177 & 1.047 & 1.047 & 0.965 \\
2.809 & 1.069 & 1.069 & 1.014 \\
4.670 & 1.124 & 1.124 & 1.092 \\
12.909 & 1.416 & 1.416 & 1.184 \\
\bottomrule
\end{tabular}
\caption{FCT/CDM ratio of the 3D power spectrum at $z=3$. ``Baryons''
counts the converted particles together with the gas; ``gas only'' counts
what is still gas. The three columns are the whole argument.}
\label{tab:3d}
\end{table}

The matter power spectra of the two runs are identical on large scales to
one part in a thousand. So is the baryon field, when the converted
particles are counted. The gas field alone is down by $19\%$.

This kills the projection as an explanation and it kills any additive
model. If the FCT excess power were entering the 1D spectrum additively
through the projection integral,
\begin{equation}
\Pod(k_\parallel) = \frac{1}{2\pi}\int_{k_\parallel}^{\infty} k\,P_{\rm 3D}(k)\,{\rm d}k ,
\end{equation}
a constant excess $C$ would contribute $(C/4\pi)(K^2 - k_\parallel^2)$,
which is positive for every $k_\parallel$ inside the range where the model
is defined and diverges where it changes sign. Our ratio does neither: it
saturates at $0.78$ and stays there. A multiplicative term is the only
thing that can sit below unity without running away, and Table~\ref{tab:3d}
is that term, measured.

The correct structure is therefore
\begin{equation}
P^{\rm gas,FCT}_{\rm 3D}(k) = b^2\,P^{\rm gas,CDM}_{\rm 3D}(k) + \Delta P_{\rm 3D}(k) ,
\end{equation}
and because the projection is linear, the 1D ratio is
\begin{equation}
\frac{\Pod^{\rm FCT}}{\Pod^{\rm CDM}}(k_\parallel)
= b^2 + \frac{\Delta \Pod(k_\parallel)}{\Pod^{\rm CDM}(k_\parallel)} .
\end{equation}
At low $k_\parallel$ the second term is negligible and the ratio goes to
$b^2 \simeq 0.81$. At high $k_\parallel$ the CDM power falls, the second
term grows, and the ratio crosses unity where the excess reaches
$1 - b^2 = 19\%$ of the CDM power. This has one free parameter fewer than
it looks: $b^2$ is not fitted, it is read off Table~\ref{tab:3d}.

Note also that $b^2$ has to sit inside the projection integral, not outside
it: only $35\%$ of the integral at $k_\parallel = 0.21\,{\rm Mpc}^{-1}$
comes from $k < 2k_\parallel$, so a scale-dependent $b^2(k)$ cannot be
pulled out. Doing the weighted average properly reproduces the measured 1D
ratio to $1\%$.

\section{Where the gas was removed matters more than how much}

Table~\ref{tab:3d} says the gas is a worse tracer. The next question is
whether that is simply because there is less of it.

It is not. Two controls settle it.

\textbf{Random removal.} Deleting gas particles at random, everywhere,
changes nothing: at a removed fraction of $0.40$ the large-scale power
ratio is $0.9999$. FCT has removed $0.404$ of the gas mass and sits at
$0.808$. It is nowhere near the null curve.

\textbf{Removal by density threshold.} Cutting CDM gas above a density
threshold chosen so that the remaining mass equals FCT's gas mass costs
$90\%$ of the large-scale power --- the ratio drops to about $0.10$. That
is the opposite extreme.

So FCT sits between two limits: random removal costs nothing, cutting the
filaments costs everything, and QLA costs $19\%$. Its removal is
concentrated enough to matter and diffuse enough not to destroy the field.
That is the signature of gas being locked into compact, highly biased
objects rather than being taken out of the diffuse IGM.

\begin{figure}[h]
\centering
\includegraphics[width=0.85\textwidth]{../figures/curve_null.png}
\caption{The removal curve. Random removal (null control) against removal
above a density threshold, with the FCT point marked. \pending{regenerate
with \code{legacy/removal\_curve\_pk.py} through the new stage}}
\end{figure}

\section{What it looks like on a single sightline}
\label{sec:t8}

The two boxes have identical initial conditions and the sightlines are
drawn from the same seed, so line $i$ passes through the same structure in
both runs. That makes a pixel-by-pixel comparison meaningful, and it is
worth doing before trusting any ensemble average.

Figure~\ref{fig:t8} shows three sightlines. The absorption troughs are in
the same places in both runs --- they have to be, the density field is the
same --- but in FCT they are systematically shallower where CDM is
strongest. The peak census makes it quantitative: over pixels with
$\tau_{\rm CDM} > 2$, the mean FCT optical depth at the same pixel is
\pending{run \code{tests/t8\_single\_los.py}}.

The single-line power spectra behave the same way as the ensemble: power
lost at low $k$, roughly tracking CDM at high $k$. The effect is not driven
by a handful of outlier sightlines.

\begin{figure}[h]
\centering
\includegraphics[width=\textwidth]{../figures/t8_p1d.png}
\caption{Individual sightlines (thin) against the ensemble (thick), and
their ratios. \pending{generate}}
\label{fig:t8}
\end{figure}

\section{Configuration space}

Everything above is a Fourier-space statistic with log binning, a Nyquist
cut, and an interpolation onto a common grid. Each of those is somewhere an
artefact could be born. The flux correlation function uses the same field
with none of that machinery, so it is a control on the estimator.

A $19\%$ suppression below $k = 0.01$~\kv{} must appear as a suppression of
$\xi_F$ on separations above roughly $600\,{\rm km\,s^{-1}}$. It does:
\pending{run \code{stages/05\_xi.py}}. The Fourier result is not a product
of the binning.

\section{Putting the converted gas back}
\label{sec:t7}

The remaining question is not whether the effect is real but what it is a
property of. Table~\ref{tab:3d} already contains the answer in outline: the
baryon field, counting the converted particles, has a ratio of $1.000$. The
gas field alone has $0.808$. The entire difference is the material QLA
removed.

So we put it back. We rebuild the snapshot with the converted particles
returned to \code{PartType0}, assign them a smoothing length from the local
baryon density and a temperature from the $T$--$\rho$ relation at their own
overdensity, and re-extract the forest with the same extractor. To show the
result does not hinge on the temperature we invented for them, we repeat
the whole thing with $T_0$ scaled by $0.5$ and by $2$.

Result: \pending{run \code{tests/t7\_stars\_back.py} and re-extract}.

This is the test that decides what we are allowed to claim. If the ratio
returns towards unity, the large-scale suppression is a property of the
star-formation prescription and not of the primordial spectrum.

\section{What this means}

Three statements, ordered by how well we can defend them.

\textbf{The suppression exists and is robust.} Five independent attempts to
remove it failed, it is present before any normalisation, and it appears in
configuration space as well as in Fourier space. This is solid.

\textbf{It comes from the gas, not from the matter.} The 3D matter power of
the two runs agrees to $0.1\%$ while the gas power differs by $19\%$. The
Lyman-$\alpha$ constraint on this class of model acts through the bias of
the surviving gas, not through the matter power spectrum. This is solid,
and as far as we know it is not how the constraint is usually presented.

\textbf{Its amplitude is $19\%$.} This is weak, and we should say so. QLA
converts $50\%$ of the baryons while the observed stellar fraction at $z=3$
is a couple of percent. A run with feedback would return much of that gas
to the IGM, the bias would recover in part, and the deficit would shrink by
an amount we cannot currently estimate.

The honest form of the result is therefore not that the model is ruled out,
and not that the signal is an artefact to be corrected away. It is that we
have identified the channel through which the forest constrains these
models, and that pinning a number on it requires a run with feedback. That
is the next paper rather than a hole in this one.

\appendix

\section{Conventions}

$F = e^{-A\tau}$; $\delta_F = F/\langle F\rangle - 1$ with
$\langle F\rangle$ taken over all pixels of all sightlines, not per
sightline;
$\Pod(k) = \langle |\,\widetilde{\delta_F}\,\Delta v|^2\rangle / L$ with
$L = N_{\rm pix}\Delta v$ and $k = 2\pi\,\mathrm{rfftfreq}(N_{\rm pix},
\Delta v)$. Plotted as $k\Pod/\pi$, which is dimensionless. The
implementation is \code{common/p1d.py} and its unit tests are in
\code{tests/test\_estimator.py}.

Sightlines are drawn from a seeded generator and the same seed is used in
both runs, which is what makes the matched comparison of
Section~\ref{sec:t8} possible.

Photoionisation rate $\Gamma_{\rm HI}(z=3) = 8.2758\times10^{-13}\,
{\rm s^{-1}}$ from HM12 via TREECOOL. The UV background table is a
command-line argument, not a hard-coded choice.

\section{Bugs found and fixed along the way}

Recorded because they were expensive and because the same traps are waiting
for anyone who repeats this.

\begin{enumerate}
\item \textbf{Truncated sightlines.} An extraction range set in Mpc against
a box specified in Mpc$/h$ silently dropped the last $31.9\%$ of every
sightline. Diagnosed as exactly one gap per line, starting at
$40.279 \pm 0.053$ and ending at $58.477 \pm 0.053$.
\item \textbf{SPH deposition without an impact parameter}, overestimating
$\tau$ by $5.6\times$.
\item \textbf{Double counting in distributed output.} With
\code{distributed:1} SWIFT writes a virtual \code{.hdf5} alongside the
\code{.N.hdf5} pieces, and a naive glob reads every particle twice. Detect
it through \code{Header/Virtual}.
\item \textbf{Catastrophic cancellation} in the Tepper-Garc\'ia
approximation to the Voigt profile in the far wings.
\item \textbf{Gas particles with negative mass} in the initial conditions,
caused by the FCT boost driving the baryon density negative in voids.
\end{enumerate}

\section{Reproducing this}

\begin{verbatim}
git clone https://github.com/patricio-c/Lyman-Alpha-PBHs && cd Lyman-Alpha-PBHs
conda activate astro
bash scripts/migrate.sh /data/contrib/pad_140/pcolazo/LOS
bash scripts/verify.sh

python tests/t0_rescaling.py  --cdm cache/cache_cdm.npz \
                              --fct cache/cache_fct.npz --out figures/t0_rescaling
python stages/04_p1d.py cache/cache_cdm.npz cache/cache_fct.npz \
                              --labels CDM FCT --out figures/p1d
python stages/05_xi.py  cache/cache_cdm.npz cache/cache_fct.npz \
                              --labels CDM FCT --out figures/xi
python tests/t8_single_los.py --cdm cache/cache_cdm.npz \
                              --fct cache/cache_fct.npz --pick median --out figures/t8
\end{verbatim}

Every cache carries the git hash, the command line and a hash of the input
snapshot, so any figure in this note can be traced back to the code that
made it.

\end{document}
